\documentclass[11pt]{article}
\usepackage[margin=1in]{geometry}
\usepackage{amsmath,amssymb,amsthm,hyperref}
\newtheorem{proposition}{Proposition}
\newcommand{\supp}{\operatorname{supp}}
\newcommand{\pc}{\operatorname{popcount}}
\title{Can the target count use fewer independent slots?}
\author{Project research note}
\date{17 November 2026}
\begin{document}
\maketitle

The current conditional bound pays, in the logarithm of expected time,
for a factor base of scale $b/m$ and a count of scale
$mb/\sqrt{\log_2 b}$. Here $b=\lceil\log_2p\rceil$ and $m$ is the
number of slots. If the count could use effective arity $\sqrt m$
at the same success probability, the balance would move from
$b/(\log_2 b)^{1/4}$ to $b/(\log_2 b)^{1/3}$, by taking
$m\asymp(\log_2 b)^{1/3}$. We test this possibility using repeated
hashes, the family of independent hash shifts, and explicit blocks.
Exact counting reduces to single coefficients of powers modulo small
prime squares. For constant-probability detection at each prescribed
nonidentity target, primes of order $m$ suffice. The cost of the
second base-prime digit is still open, and the $1/3$ rate remains unproved.

Write $G=E(\mathbb F_p)$ additively. A slot $F$ is the set of
affine points with $(\alpha x+\beta\bmod M)<\tau$, where $M=2^b$,
$\alpha$ is odd, and $\tau$ is a power of two smaller than $M$.
Thus $F=-F$ and $|F|\le 2\tau$. For the project's paid width,
$\log_2\tau=b/m+O(\log b+\log m)$. The baseline and its offset
reduction are in \texttt{../short\_hash\_count/count.tex}.
Our algebra concerns the ordinary ordered count at one target.
Guarded C38 counts have extra restrictions on intermediate sums.

\section{What one bit of the count can see}
Use the integral group algebra, writing $[P][Q]=[P+Q]$.
Suppose the $m$ slots consist of masks $h_j=\sum_{P\in F_j}[P]$
repeated $m_j$ times, for $1\le j\le g$, and let $H=\prod_j h_j^{m_j}$.
For $e\ge0$ define $D_e h_j=\sum_{P\in F_j}[2^eP]$.
If $I_j$ is the set of one-bits in $m_j$, then
\[
   H\bmod 2=\prod_j\prod_{e\in I_j} D_e h_j\bmod 2 .
\]
Indeed, squaring in the commutative group algebra in characteristic
two doubles each point, and the binary expansion supplies the product.
The new arity is $s=\sum_j |I_j|$. In an odd-order group, for a
single mask and $m=2^e$, parity at $T$ is exactly the one membership
bit $1_F([2^e]^{-1}T)$, where the inverse is that of the doubling
automorphism. No table of sums is involved. Parity can certify a
nonzero integer count only on the support of this reduced product,
of size at most $(2\tau)^s$. In the one-mask power-of-two case the
price is exact for every fixed $T\ne O$. With a uniform shift, the
unique point $[2^e]^{-1}T$ passes in exactly $\tau$ of the $M$
shifts. The expected number of independent trials until odd parity
is $M/\tau=2^{(1-1/m)b+O(\log b+\log m)}$.

\begin{proposition}
For precision $q\ge1$, every target visible in $H\bmod 2^q$ is a
sum of at most $d=\min(m,s+q-1)$ terms of the form $2^e P$ with
$P\in F_j$ and $2^e\le m_j$. In particular, for every choice of seeds,
\[
 |\supp(H\bmod 2^q)|\le
 \left(1+2\tau\sum_j(1+\lfloor\log_2 m_j\rfloor)\right)^d .
\]
\end{proposition}
\begin{proof}
Group the ordered tuples by their multiplicities $n_{j,P}$.
Each class has weight $\prod_j m_j!/\prod_{j,P}n_{j,P}!$. Using
$v_2(u!)=u-\pc(u)$, its valuation is
$\sum_{j,P}\pc(n_{j,P})-s$.
If $H(T)$ is nonzero modulo $2^q$, at least one class for $T$
has valuation less than $q$. Expand each of its multiplicities in
binary. This expresses $T$ as the sum of at most $s+q-1$ terms,
also at most $m$. Pad to $d$ terms with the identity and count the
choices. Several encodings can give the same target.
\end{proof}

On the prime-order subclass let $r=|G|$, so $\log_2 r=b+O(1)$.
For a uniform nonidentity target independent of the masks, the
success probability of the nonzero-residue test is at most
\[
 \min\{1,\ 2^{-(1-d/m)b+O(m(\log b+\log m))}\}.
\]
This applies also after averaging over any seed distribution independent
of that target. For polylogarithmic $m$, even $d=m-1$ makes this
probability $2^{-b/m+O(m(\log b+\log m))}$. Thus constant visibility
requires $q\ge m-s+1$ asymptotically. If $s=O(\sqrt m)$ and $q=1$,
the retry cost for independent uniform trials is $2^{(1-o(1))b}$
as $m\to\infty$. This argument prices a test that accepts on a
nonzero residue. A zero can conceal many tuples, and other existence
tests and target-dependent seed choices have different coverage laws.

Widening the hashes also has a price. For parity and constant
visibility at uniform targets, the support bound alone requires
$s\log_2(2\tau)\ge b-O(1)$. The existing factor-base enumeration
prices $\tau$ candidates per mask. If the reduced $s$-factor
product is then priced at the existing long-coordinate rate,
the two charges have scales at least $b/s$ and
$s b/\sqrt{\log_2 b}$, respectively, up to constants and additive
lower-order terms. Optimizing that ledger gives the original
fourth-root logarithmic saving. Even granting the fast contraction
estimate for the doubled masks would not certify the proposed
third-root saving by this change of width.

\paragraph{The precise bit at which independent slots return.}
For arbitrary slot indicators $f_1,\ldots,f_m$ put
\[
 A=\sum_{S\subseteq[m]}(-1)^{m-|S|}
       \left(\sum_{j\in S} f_j\right)^m .
\]
Expansion cancels every monomial omitting a slot. The surviving
multiplicity is $m!$, so $A=m!\prod_i f_i$ in the integral group
algebra. Set $v=m-\pc(m)$. The low $v$ bits of $A(T)$ vanish, and
bit $v$ (numbering from zero) equals the parity of the independent
target count. Thus a routine for powers of \emph{sums} of hash masks
at precision $v+1$ suffices for that parity, at $2^m$ calls. For
polylogarithmic $m$ the extra $m$ in the log of the call count is
lower order in the rates considered here. Such sums are weighted
functions and need not be individual dyadic masks. For exact counts
bounded by $U$, precision $v+\lceil\log_2(U+1)\rceil$ suffices by
dividing out $2^v$ and inverting the odd part of $m!$. These identities
locate the necessary precision without providing a fast power routine.
For the classical polarization identity, see
\href{https://arxiv.org/abs/1309.1275}
{E. G. F. Thomas, A polarization identity for multilinear maps}.

\section{Exact reconstruction using only prime squares}
There is a way around the increasing binary precision. For a prime
$l\ge m$ define the additive map $\Phi_l([P])=[lP]$, extended
$\mathbb Z$-linearly. On the integral group algebra set
\[
        D_l(h)=\frac{h^l-\Phi_l(h)}{l}\bmod l .
\]
The numerator is coefficientwise divisible by $l$, by the Frobenius
identity and Fermat's little theorem on integer coefficients.
Evaluating $D_l(h)(T)$ requires only the numerator modulo $l^2$.
For $r$ prime and $l\ne r$, $\Phi_l(h)(T)$ is one evaluation of $h$
at $[l]^{-1}T$. Its cost is polynomial for the masks below on the
supplied-order input. This leaves the coefficient of $h^l$ modulo
$l^2$ to compute.

For nonnegative coefficients there is an integer interpretation. Give
point $P$ exactly $h(P)$ colors. Cyclic rotation acts on length-$l$
words of colored points with sum $T$. Its fixed words are constant
and there are exactly $\Phi_l(h)(T)$ of them. Every other orbit has
size $l$, so $D_l(h)(T)$ is the number of these orbits modulo $l$.
Listing the orbits explicitly can be expensive even for one genuine
unweighted hash mask of size $a$. At least $a^l-a^{l-1}-a$ words
are nonconstant and sum to a nonidentity target, since at most
$a^{l-1}$ sum to the identity. Some nonidentity target has at least
$(a^l-a^{l-1}-a)/(l(r-1))$ orbits. At the paid threshold a shift
with $a\ge r^{1/m}$ exists by the average $(r-1)\tau/M$. For any
queried prime $l\ge2m$ this example has $\Omega(r/l)$ orbits at
one target. The digit is one small residue even in this case.

\begin{proposition}
For a known bound $0\le H(T)=(f_1*\cdots*f_m)(T)\le U$, the
exact integer $H(T)$ reduces deterministically to
$2^m\operatorname{poly}(m,\log(U+2))$ queries for single coefficients
$D_l(h)(T)$. Each $h$ is a sum of some original slot indicators and
a nonnegative integer multiple of $[O]$. The primes have size
polynomial in $m+\log(U+2)$.
\end{proposition}
\begin{proof}
For each prime $l\ge m$, pad the list of $m$ inputs with $l-m$
copies of $[O]$. Apply the integral polarization identity at arity
$l$. The alternating sum of all $\Phi_l$ terms vanishes by additivity
since $l\ge2$. Divide by $l$ before reducing modulo $l$.
Group identical padding inputs by the number $t$ selected to obtain
\[
 (l-1)!\, H(T)\equiv
 \sum_{S\subseteq[m]}\ \sum_{t=0}^{l-m}
 (-1)^{l-|S|-t}\binom{l-m}{t}
 D_l\left(t[O]+\sum_{j\in S} f_j\right)(T) \pmod l .
\]
The factorial $(l-1)!$ is invertible modulo $l$. Equivalently,
query each raw coefficient of $h^l$ modulo $l^2$, accumulate the
signed sum modulo $l^2$, and divide its canonical residue by $l$.
The $\Phi_l$ terms cancel across the entire sum. This version never
evaluates a pushforward and applies even if $l$ divides $|G|$.
There are at most $2^m(l-m+1)$ queries regardless of the size
of the padding. Use distinct primes with product exceeding $U$,
then recover the unique integer in $[0,U]$ by the Chinese remainder
theorem. For example, $L=\max(1,\lceil\log_2(U+1)\rceil)$ primes
greater than $m$ suffice. A sieve up to a polynomial in $L+m$ finds
them. Indeed $\binom{2N}{N}\ge 4^N/(2N+1)$, and its valuation at
each prime $a$ is at most $\lfloor\log_a(2N)\rfloor$, by subtracting
the factorial floor sums. Thus $\binom{2N}{N}\le(2N)^{\pi(2N)}$,
giving sufficiently many primes by taking $N$ of order $(L+m+2)^2$.
The binomial weights, arithmetic modulo $l$, and integer reconstruction
have polynomial overhead per query. The sum mask is specified by $S,t$
and the seeds. The reduction never materializes a partial-sum support.
\end{proof}

In the project, $U=(2\tau)^{m-1}$ suffices since the final group
element is determined by the other $m-1$. All requested powers thus
have modulus polynomial in $b,m$ (and word length $O(\log(b+m))$).
At polylogarithmic arity the extra $m+O(\log b)$ in the logarithm of
the query count is lower order than $b/(\log_2 b)^{1/3}$.
The apparent reduction to an easy Frobenius map fails specifically
in the next digit. For example,
\[
 D_l(A+B)-D_l(A)-D_l(B)
   =\sum_{j=1}^{l-1}\frac{(-1)^{j-1}}{j} A^j B^{l-j}\pmod l
\]
by $\binom{l}{j}/l=\binom{l-1}{j-1}/j$. The required second digit
contains interactions between the masks even for one target. A fast
algorithm for this digit with suitably uniform dependence on $l$
would advance exact counting. Calling the existing contractor on
$l$ factors does not establish such an algorithm.

\paragraph{Linear-size primes already suffice for constant detection probability.}
The preceding deterministic reconstruction also handles exceptional
seeds with many tuples. For the project's distribution one can
enforce $m<l\le16m$ when the desired output is a certificate that
the ordinary count is positive. Here let $N=|G|$ with $N\ge30m$ as in
the supplied coverage lemma. This argument allows composite $N$ as
the raw-power implementation requires no inverse of $l$ on $G$.
For each fixed $T\ne O$, consider
all ordered affine tuples summing to $T$. There are at most
$N^{m-1}$, and each survives the independent shifts with probability
$(\tau/M)^m$. Thus
\[
  \mathbb E H(T)\le (N\tau/M)^m/N < 4^m .
\]
The strict bound uses the smallest paid dyadic threshold in the
regime $2N^{1/m}\ge20m$. This holds eventually for polylogarithmic
$m$ as $b$ grows. The inherited targetwise coverage bound gives
$\Pr(H(T)>0)>2/5$ because every guarded tuple is ordinary.
Take primes with product $J>10\cdot4^m$. A nonzero residue provides
a true certificate, and every $0<H(T)<J$ is detected. Markov's
inequality therefore gives detection probability greater than
$2/5-1/10=3/10$ for each fixed nonidentity target.

For sufficiently large $m$, the primes in $(m,16m]$ have more
than enough product. In fact $\binom{16m}{8m}\ge 2^{16m}/(16m+1)$.
All prime factors after their first copy contribute at most
$(16m)^{\sqrt{16m}}$, using the valuation bound above. Conversely
the product of primes through $m$ is at most $2^{4m}$, by rounding
up to a power of two and bounding each dyadic interval's product
by its central binomial coefficient. Hence the logarithm of the
interval product is at least $12m-O(\sqrt m\log m)$ in base two.
The total number of digit queries is $O(2^m m^2)$. This sharpens the
missing algorithm to prime powers with degree only linear in the
original arity for this distributional nonemptiness test. The test
certifies ordinary tuples. The guaranteed detection probability and
the cost of evaluating a digit are distinct parts of the accounting.
For $\ell>m$, the later calculation in
\texttt{../second\_digit/digit.tex} takes the padding differences
before querying any digits. They cancel the unpadded digit and leave
a degree-$m$ mixed coefficient. A fast individual digit routine is
sufficient for this reduction, and the combined query admits that
additional simplification.

\section{Independent seeds at one fixed target}
Assume $G$ has odd prime order $r\ge5$, let $n=(r-1)/2$ and
let $X$ be the $n$ affine abscissas. Keep every multiplier fixed and
vary the shifts of the actual hash family. Consider a split of the slots
into two nonempty groups. A separated evaluation with $q$ scalar terms
has the form $\sum_{\nu=1}^q L_\nu(\text{left shifts})
R_\nu(\text{right shifts})$. At a fixed target its matrix of values
has rank at most $q$. We bound this rank with no free target index.

Let $c$ be the number of residue classes modulo $\tau$ fully contained
in $X$ when these classes are taken inside $\mathbb Z/M$. Define
$\delta=\max(c-1,0)$. For each fixed odd $\alpha$, the span of the
shift masks has dimension $n-\delta$ over any field. Indeed, extend
an annihilating vector on $\alpha X$ by zero to the full cycle.
Every window of length $\tau$ sums to zero. Subtract adjacent window
conditions to obtain $\tau$-periodicity. Hence the vector is constant
on each complete class, vanishes on the others and has sum zero over
one period. Conversely these conditions suffice. Multiplication by
$\alpha$ permutes the residue classes and preserves completeness.
This is a calculation on abscissas, each representing one signed
point pair, with no division by two in characteristic two. If
$M-p\ge\tau$ then $c=0$ because the omitted canonical tail meets
every residue class.

\paragraph{Rational arithmetic.}
Fix a target $T$ and all shifts except one in each group. Choose the
fixed shifts so their masks are nonempty. Their convolution $k$
(the identity for the empty collection) is invertible over
$\mathbb Q$. Each mask polynomial in a cyclic coordinate of $G$ has
nonzero value at one and at every nontrivial $r$th root of unity,
since the latter have minimal polynomial $1+z+\cdots+z^{r-1}$ and
the indicator is nonconstant. No cyclic coordinates need to be
computed by an algorithm for this proof.
On the $(n+1)$-dimensional space of even group functions the form
$(f,g)\mapsto(f*g*k)(T)$ is nondegenerate. In Fourier coordinates
the nontrivial characters pair by inversion, with target multiplier
$\zeta^{iT}+\zeta^{-iT}\ne0$ for $i\ne0$ because $r$ is odd.
The trivial component is nonzero as well. Restricting each side to a
space of codimension $1+\delta$ gives the lower bound
\[
       \operatorname{rank}_{\mathbb Q} \ge n-1-2\delta .
\]
It holds for every target and every nonempty choice of the fixed masks.
Since $c\le\tau$, the bound has size $2^{b-O(1)}$ at paid polylogarithmic
arities tending to infinity. The asymptotic statements here are bounds
on the prime-order instances as their size increases.

When $\delta=0$ and each group has at least two slots, the rank is
exactly $n+1$. The one-slot span is all even functions vanishing at
the identity, with basis $e_P=[P]+[-P]$, one per sign pair. The
two-slot products span the full even space. Indeed, if a functional
$l$ annihilates them, $e_P^2=e_{2P}+2[O]$ forces
$l(e_Q)=-2l([O])$ for every $Q\ne O$. For $P\ne\pm Q$,
$e_Pe_Q=e_{P+Q}+e_{P-Q}$ then forces $4l([O])=0$. Further slots
can each be contracted to one virtual $e_P$ by linear combinations
of shifts. This element is invertible over $\mathbb Q$ by the same
root argument. Pairing the two full spans gives the claim.

\paragraph{Parity.}
At every $T\ne O$ the rank of the full matrix over $\mathbb F_2$
is at least $n-2\delta$ and at most $n$. To see this, the subspace
$V$ of even functions with zero at $O$ equals the even functions
with zero augmentation (sum of coefficients) in this characteristic.
It is closed under convolution and has dimension $n$. Over the
algebraic closure, the Fourier transform identifies it with the
$n$ independent nontrivial pairs of character coordinates. This
transform exists because $r$ is odd. The pairing at $T\ne O$ has
diagonal factors $\zeta^{iT}+\zeta^{-iT}\ne0$, so it is
nondegenerate on $V$. At $O$ it vanishes.

For each extra slot take a nonzero difference of adjacent windows.
Such a difference exists by the positive dimension of the slot span
and the fact that each canonical abscissa is retained by some shifts
and omitted by others. It is a signed pair $e_P$, or a difference
of two distinct signed pairs $e_P-e_Q$, up to sign. Both act
invertibly on $V$. In characteristic two $z+z^{-1}=0$ for an
$r$th root only if $z=1$, and $z+z^{-1}=w+w^{-1}$ implies
$z=w$ or $z=w^{-1}$. A nontrivial character of a prime-order group
is faithful. Therefore none of the $n$ multipliers vanishes.
Contract these extra shift axes by the corresponding two coefficients
on each side, leaving one free slot on each side. This operation
is a linear row and column transformation of the full matrix.
The two remaining spaces in $V$ have codimension $\delta$, giving
the stated lower bound. Ranks are unchanged by extending the field.
If actual fixed shifts are required, expand the product of window
differences. Rank subadditivity shows some single choice of the
other shifts has rank at least $(n-2\delta)/2^{m-2}$.

Thus a short sum of separated scalar products across independent
hash seeds fails even at one prescribed nonidentity target and even
with parity output. This is an exact-rank result. Approximation,
expected running time over seeds, and computations that read both
groups jointly are outside its lower bound.

\section{Charging materialized blocks}
For completeness consider grouping the slots into $a$ blocks and
handing full supports to an outer contraction. At the original
independent-hash coverage threshold, with probability at least $1/4$,
the whole support contains at least $(r-1)/5$ targets. This follows
by averaging the supplied targetwise coverage probability $>2/5$
over the nonidentity targets, as in \texttt{../c38\_arithmetic\_join/join.tex}.
On that same event, simultaneously for every partition (including
a seed-dependent one), the product of the $a$ block support sizes
is at least $(r-1)/5$. Consequently the largest table has log size
at least $(b-O(1))/a$.

If the proposed outer step is priced using the current $a$-slot
long-coordinate estimate of order $ab/\sqrt{\log_2 b}$, then even
granting that estimate for table inputs and ignoring their interface
cost, the two lines of the budget obey
\[
 \max\left\{\frac{b-O(1)}a,\frac{a b}{\sqrt{\log_2 b}}\right\}
 \ \ge\ (1-o(1))\frac{b}{(\log_2 b)^{1/4}} .
\]
The second line is a charge in this proposed analysis. It is not a
lower bound on the runtime of all outer contractions. The existing
$26.2933775156+o(1)$ time coefficient and its conditional status are
unchanged by these tests. The open algorithmic question is whether
one can evaluate the full mixed coefficient using joint access to
the seeds, below the current arity cost, with sufficient visibility
and witness recovery under the paid target distribution.

\paragraph{Finite checks.}
Run \texttt{python -B research/slot\_arity/verify.py} from the
project root. The order-19 and order-31 curves check rational
one-target lower bounds with fixed background masks and balanced-cut
parity ranks at every target. For rational balanced ranks the checker
certifies a lower bound by a minor modulo 1009, matching the
characteristic-zero dimension upper bound. Tuple convolution checks
the repeated-mask bit support and polarization. On the order-19
fixture eight equal four-point slots reach all 19 targets as ordinary
counts, of which four have odd count. The prime-square reconstruction
recovers ordinary counts as large as 122 and 568 on the two curves.
The reference digit oracle computes literal powers on the tiny group.
An exhaustive three-slot shift average verifies the first-moment
identity at every target over 4096 independent shift choices.
An additional curve
$y^2=x^3+9x+1$ over $\mathbb F_{1031}$ has order 1019 and, at the
original eight-slot coverage threshold $\tau=512$, all 2048 shifts
at multiplier seven give one-slot rank 509 modulo two. Its balanced
fixed-target ranks of 510 over the rationals and 509 modulo two at
nonidentity targets follow from the proofs. The large balanced
matrix is never constructed. These finite checks measure algebraic
ranks and coefficients, not solver runtimes.
\end{document}
